% Claim proof file for row claim_3_18 -- "MarginalizationAndIntervention".
%
% Proof lifted verbatim from lecture-notes/lecture_notes/graphs.tex lines
% 1131--1164 (the LN's `\Claude{...}` wrapper stripped).
%
% Compiles standalone via the `subfiles` package.

\documentclass[main]{subfiles}

\begin{document}
% ref: claim_3_18 (proof, with statement restated for readability)
% title: MarginalizationAndIntervention
\def\rowref{claim\_3\_18}
\phantomsection\label{claim_3_18}
%
% Cross-references: see claim_<ref>_statement_<title>.tex in this folder
% for the corresponding statement.

% --- statement (restated) ---------------------------------------------------
\begin{Lem}[Marginalization and intervention commute]
      Let $G=(J,V,E,L)$ be a CDMG and $W_1 \ins J \cup V$ and $W_2 \ins V$ two disjoint subsets of nodes from $G$.
      Then we have:
      \[ \lp G_{\doit(W_1)} \rp^{\sm W_2} = \lp G^{\sm W_2} \rp_{\doit(W_1)}. \]
     % This shows that the following notations are unambiguous:
     % \[G(V\sm (W_1 \cup W_2)|\doit(W_1 \cup J))=G_{\doit(W_1)}^{\sm W_2}.\]
      A similar statement holds for marginalizations and adding intervention nodes, and also for marginalizations and node-splitting interventions.
\end{Lem}

% --- proof ------------------------------------------------------------------
\begin{proof}
We verify that all four components of $\lp G_{\doit(W_1)} \rp^{\sm W_2}$ and $\lp G^{\sm W_2} \rp_{\doit(W_1)}$ agree.

\medskip\noindent\emph{Input nodes.}
Both sides have $J \cup W_1$: marginalization does not change input nodes, and hard intervention adds $W_1$ to them.

\medskip\noindent\emph{Output nodes.}
Left: $(V \sm W_1) \sm W_2 = V \sm (W_1 \cup W_2)$. Right: $(V \sm W_2) \sm W_1 = V \sm (W_1 \cup W_2)$.

\medskip\noindent\emph{Directed edges.}
Let $\ul{v}, \ol{v} \in (J \cup W_1) \cup (V \sm (W_1 \cup W_2))$.

Left: $\ul{v} \tuh \ol{v} \in (E_{\doit(W_1)})^{\sm W_2}$ iff there is a directed walk
$\ul{v} \tuh a_1 \tuh \cdots \tuh a_{n-1} \tuh \ol{v}$ in $G_{\doit(W_1)}$ with all $a_k \in W_2$.
Since $E_{\doit(W_1)} = E \sm \{v \tuh w \mid w \in W_1\}$ and all $a_k \in W_2$ satisfy $a_k \notin W_1$ (as $W_1 \cap W_2 = \emptyset$), such a walk exists iff there is a directed walk in $G$ from $\ul{v}$ to $\ol{v}$ with intermediate nodes in $W_2$ and no edge pointing into $W_1$.

Right: $\ul{v} \tuh \ol{v} \in (E^{\sm W_2})_{\doit(W_1)}$ iff $\ul{v} \tuh \ol{v} \in E^{\sm W_2}$ and $\ol{v} \notin W_1$.
The latter holds since $\ol{v} \in V \sm (W_1 \cup W_2)$.
And $\ul{v} \tuh \ol{v} \in E^{\sm W_2}$ iff there is a directed walk in $G$ from $\ul{v}$ to $\ol{v}$ with intermediate nodes in $W_2$.
Since all $a_k \in W_2 \ins V \sm W_1$, no edge in such a walk points into $W_1$, and $\ol{v} \notin W_1$.
Hence both conditions are equivalent.

\medskip\noindent\emph{Bidirected edges.}
Let $\ul{v}, \ol{v} \in V \sm (W_1 \cup W_2)$, $\ul{v} \neq \ol{v}$.

Left: $\ul{v} \huh \ol{v} \in (L_{\doit(W_1)})^{\sm W_2}$ iff there is a bifurcation in $G_{\doit(W_1)}$ between $\ul{v}$ and $\ol{v}$ with intermediate nodes in $W_2$.
Since $L_{\doit(W_1)} = L \sm \{v \huh w \mid w \in W_1\}$ and $W_2 \cap W_1 = \emptyset$, such a bifurcation uses only edges not touching $W_1$, so it is also a bifurcation in $G$ with intermediate nodes in $W_2$.

Right: $\ul{v} \huh \ol{v} \in (L^{\sm W_2})_{\doit(W_1)}$ iff $\ul{v} \huh \ol{v} \in L^{\sm W_2}$ and $\ul{v}, \ol{v} \notin W_1$.
The latter holds. And $\ul{v} \huh \ol{v} \in L^{\sm W_2}$ iff there is a bifurcation in $G$ between $\ul{v}$ and $\ol{v}$ with intermediate nodes in $W_2$.
Since $W_2 \cap W_1 = \emptyset$, both sides reduce to: there exists a bifurcation in $G$ between $\ul{v}$ and $\ol{v}$ with intermediate nodes in $W_2$.
\end{proof}

\end{document}
